\begin{trapbox}

	\begin{description}[style=nextline, leftmargin=2.8em, labelsep=0.5em, itemsep=0.35em]
		\item[\textbf{\textcolor{CTrap}{易错点一（\(R\) 运算失误）}}]
		      错误地将
		      \[
			      R
		      \]
		      计算为
		      \[
			      \sqrt{aA^2+bB^2}
		      \]
		      或
		      \[
			      \sqrt{a^2A+b^2B},
		      \]
		      忽略了 \(a,A,b,B\) 都应以平方形式出现。

		\item[\textbf{\textcolor{CTrap}{正确理解（公式核对）：}}]
		      \[
			      R=\sqrt{a^2A^2+b^2B^2}.
		      \]
		      注意 \(a\) 与 \(A\) 各自平方后再相乘，\(b\) 与 \(B\) 也一样。

		\item[\textbf{\textcolor{CTrap}{错因分析（顺序混淆）：}}]
		      学生容易混淆平方的顺序，或误认为
		      \[
			      R=aA+bB.
		      \]
		      建议写成
		      \[
			      R=\sqrt{(aA)^2+(bB)^2},
		      \]
		      这样更不容易漏平方。
	\end{description}

	\medskip

	\begin{description}[style=nextline, leftmargin=2.8em, labelsep=0.5em, itemsep=0.35em]
		\item[\textbf{\textcolor{CTrap}{易错点二（绝对值忽略）}}]
		      直接使用 \(C\) 代入公式，忘记取 \(|C|\)，导致最值符号错误。

		\item[\textbf{\textcolor{CTrap}{正确理解（符号处理）：}}]
		      公式中出现的是
		      \[
			      |C|,
		      \]
		      无论 \(C\) 的正负，都要取绝对值。

		\item[\textbf{\textcolor{CTrap}{错因分析（惯性思维）：}}]
		      学生习惯于使用直线方程系数原样代入，忽略距离公式中分子本身带绝对值。
	\end{description}

	\medskip

	\begin{description}[style=nextline, leftmargin=2.8em, labelsep=0.5em, itemsep=0.35em]
		\item[\textbf{\textcolor{CTrap}{易错点三（相交时最小值）}}]
		      当 \(|C|\le R\) 时，错误地使用
		      \[
			      \frac{|C|-R}{\sqrt{A^2+B^2}}
		      \]
		      作为最小值，从而得到负数。

		\item[\textbf{\textcolor{CTrap}{正确理解（取零处理）：}}]
		      此时最小值为
		      \[
			      0,
		      \]
		      因为直线与椭圆有公共点，椭圆上存在点到直线的距离为零。

		\item[\textbf{\textcolor{CTrap}{错因分析（未理解几何意义）：}}]
		      学生套用公式时，没有结合几何意义判断是否有交点，导致直接计算
		      \[
			      \frac{|C|-R}{\sqrt{A^2+B^2}},
		      \]
		      得到负数后未取 \(0\)。
	\end{description}

	\medskip

	\begin{description}[style=nextline, leftmargin=2.8em, labelsep=0.5em, itemsep=0.35em]
		\item[\textbf{\textcolor{CTrap}{易错点四（把取等点误认为标准顶点）}}]
		      认为最大距离或最小距离一定在椭圆的四个标准顶点
		      \[
			      (\pm a,0),\quad (0,\pm b)
		      \]
		      处取得。

		\item[\textbf{\textcolor{CTrap}{正确理解（法向支撑点）：}}]
		      一般情况下，取等点不是标准顶点，而是与直线法向量相关的法向支撑点：
		      \[
			      P_+=\left(\frac{a^2A}{R},\frac{b^2B}{R}\right),
			      \qquad
			      P_-=-P_+.
		      \]

		\item[\textbf{\textcolor{CTrap}{错因分析（几何概念混淆）：}}]
		      “顶点”是椭圆坐标轴方向上的特殊点，而本题的距离最值由直线的法向量方向决定。除非直线法向量恰好与坐标轴方向一致，否则取等点通常不是标准顶点。
	\end{description}

	\medskip

	\begin{description}[style=nextline, leftmargin=2.8em, labelsep=0.5em, itemsep=0.35em]
		\item[\textbf{\textcolor{CTrap}{易错点五（把两条切线的较小距离恒当作最小值）}}]
		      使用平行切线法后，误以为无论原直线在什么位置，都有
		      \[
			      d_{\min}=\min\left\{\frac{|C-R|}{\sqrt{A^2+B^2}},
			      \frac{|C+R|}{\sqrt{A^2+B^2}}\right\}.
		      \]

		\item[\textbf{\textcolor{CTrap}{正确理解（先判断原直线的位置）：}}]
		      两条支撑切线为 \(Ax+By=\pm R\)。
		      当 \(|C|>R\) 时，原直线位于带状区域外，上式才是正确的最小距离。
		      当 \(|C|\le R\) 时，原直线与椭圆有交点，必须取
		      \[
			      d_{\min}=0.
		      \]
		      但是最大距离始终可以取两条支撑切线距离中较大的一个。

		\item[\textbf{\textcolor{CTrap}{错因分析（支撑切点与实际交点混淆）：}}]
		      平行切点控制的是沿已知法向方向的最外边界；
		      当原直线已经切到或穿过椭圆时，距离为零的点是切点或交点，
		      不一定是前面求出的两个法向支撑点。
	\end{description}

\end{trapbox}